% ===================================================================
%  TAVOLA N. 3 — Infanzia e giovinezza.
%  Traduction 2026 d'apres texto/io, controlee sur texto/fr et
%  texto/en. Consigne : voir texto/it/10-tavola-01.tex.
%
%  Deux scenes, et l'ido subdivise beaucoup plus que le francais :
%  huit des intertitres n'ont pas de vis-a-vis francais. L'italien
%  suit l'ido.
%
%  LE TABLEAU 3 ALLEMAND AVAIT MONTRE UNE LANGUE QUI A PERDU SES NOMS
%  DE MOIS. CELLE-CI EST LA LANGUE QUI LES A DONNES.
%
%  L'allemand avait eu Lenzing, Ostermond, Wonnemond, Brachet, Heuert,
%  Ernting, et deux campagnes de puristes n'ont pas suffi a les
%  ramener contre März, April, Mai, Juni, Juli, August. L'italien
%  n'avait rien a perdre : marzo, aprile, maggio, giugno, luglio,
%  agosto SONT les mots latins, sans une lettre de changee, et c'est
%  d'ici qu'ils sont partis pour toute l'Europe.
%
%  ET DEUX DES DOUZE SONT DES NOMS D'HOMMES QUI ONT REGNE. « Luglio »
%  est Iulius Caesar, « agosto » est Augustus. Or gravuri/korekti.json
%  met justement « luglio » en gras au bloc c2-01-1, pour corriger une
%  coquille de l'atelier ido : LE SEUL MOT DE LA PAGE QUE LA TABLE DES
%  CORRECTIONS TOUCHE EST LE NOM DE CESAR. Il est en gras dans les
%  trente-quatre colonnes, et il n'y a que celle-ci pour l'entendre.
%
%  ET LA GIOSTRA DE L'ALINEA 12 CONFIRME LE SIXIEME MECANISME DES LA
%  PAGE SUIVANTE. Le tableau 2 venait de trouver que le football
%  italien ne s'appelle ni footbal ni balle-au-pied mais
%  calcio, du nom d'un jeu florentin du XVIe siecle : le nom
%  n'est ni emprunte ni forge ni calque, il est DEPLACE d'une chose
%  ancienne a une chose neuve qui lui ressemble.
%
%  Le manege de la foire s'appelle giostra, et « giostra »
%  est LA JOUTE, le tournoi. Ce n'est pas une image : les premiers
%  maneges etaient des machines a s'exercer a la bague, ou le
%  cavalier tournait pour enfiler un anneau au bout de sa lance. La
%  chose a change, le nom est reste, et il designe aujourd'hui un
%  divertissement d'enfants. L'allemand dit « Karussell », mot
%  francais venu de l'italien « carosello » — un autre exercice
%  equestre. LES DEUX MOTS QUE L'EUROPE EMPLOIE POUR LE MANEGE SONT
%  DEUX NOMS DE TOURNOI, ET TOUS DEUX SONT ITALIENS.
%
%  Un troisieme du meme genre, a l'alinea 13 : l'arena (44) ou l'on
%  combat le taureau. « Arena » est le
%  latin pour LE SABLE — celui qu'on repandait sur le sol de
%  l'amphitheatre pour boire le sang. L'italien nomme le batiment par
%  son plancher, et le nom a fait le tour du monde sans que personne
%  y pense.
%
%  TROIS FOIS SUR DEUX PAGES, DONC. Le deplacement n'est pas un
%  accident de l'italien : c'est sa maniere ordinaire. Une langue qui
%  a trois mille ans de mots derriere elle n'a pas besoin d'en
%  fabriquer — elle en a d'avance, et il lui suffit de les remployer.
% ===================================================================

%%K t03-apar-1 apar
\VUsaut{3.0mm}
\VUtitre{15.0pt}{60}{TAVOLA N. 3}
\VUsaut{1.6mm}
\VUtitre{11.4pt}{0}{Infanzia e giovinezza.}
\VUsaut{3.4mm}

%%K t03-apar-2 apar
(Le tavole 3 e 4 sono disposte in modo da presentare, in quattro
\VUgras{scene di famiglia}, il vocabolario essenziale del \VUgras{tempo},
dell'\VUgras{età}, della \VUgras{famiglia}, del \VUgras{vestiario} e della
\VUgras{salute}.)
\VUsaut{4.0mm}

%%K t03-tit-1 sub
\VUcentre{12.0pt}{0}{\textit{Prima scena.}}
\VUsaut{3.0mm}

%%K t03-tit-2 sub
\VUpk{10.2pt}{40}{Un battesimo in paese.}
\VUsaut{2.4mm}

%%K t03-tit-3 sub
\VUpk{10.2pt}{40}{Primavera.}
\VUsaut{3.0mm}

%%K t03-c1-01-1 p
1. --- Siamo in \VUgras{primavera}, nel mese di \VUgras{marzo}, \VUgras{aprile} o
\VUgras{maggio}, in un giorno della \VUgras{settimana} --- \VUgras{lunedì},
\VUgras{martedì} o \VUgras{mercoledì}. Sono le \VUgras{dieci} del mattino.

%%K t03-c1-02-1 p
2. --- Il \VUgras{tempo} è bello, l'\VUgras{aria} è limpida. Il sole splende. Qualche
piccola \VUgras{nuvola} \textsuperscript{(2)} sale nel \VUgras{cielo}
\textsuperscript{(1)} azzurro.

%%K t03-c1-03-1 p
3. --- Gli \VUgras{alberi} non hanno quasi \VUgras{foglie}. Gli snelli
\VUgras{pioppi} \textsuperscript{(3)} e l'alto \VUgras{platano}
\textsuperscript{(17)} portano ancora soltanto gemme.

%%K t03-c1-04-1 p
4. --- Le \VUgras{rondini} \textsuperscript{(4)} sono tornate e volteggiano con grida
allegre intorno alla \VUgras{guglia} \textsuperscript{(7)} del \VUgras{campanile}
\textsuperscript{(5)} che sovrasta la \VUgras{chiesa} \textsuperscript{(6)} del paese.

%%K t03-tit-4 sub
\VUsaut{3.4mm}
\VUpk{10.2pt}{40}{La chiesa.}
\VUsaut{3.0mm}

%%K t03-c1-05-1 p
5. --- Questa chiesa è molto semplice, con il suo grande \VUgras{portale}
\textsuperscript{(13)} e il suo grande \VUgras{orologio} \textsuperscript{(8)}. La
\VUgras{lancetta} \textsuperscript{(9)} piccola è sulla X: segna le \VUgras{ore}. La
\VUgras{grande} \textsuperscript{(10)} è sulla XII (mezzogiorno); segna i
\VUgras{minuti}.

%%K t03-c1-06-1 p
6. --- Nel campanile le \VUgras{campane} \textsuperscript{(11)} suonano allegramente.
In cima al tetto c'è una piccola \VUgras{croce} \textsuperscript{(12)} di pietra.

%%K t03-c1-07-1 p
7. --- La chiesa non ha solo un grande \VUgras{portale} \textsuperscript{(13)}, ha
anche una porticina laterale. Da questa porta il \VUgras{parroco}
\textsuperscript{(15)}, con la sua \VUgras{tonaca}, esce dalla \VUgras{sacrestia}
\textsuperscript{(14)} e si avvia verso la \VUgras{canonica} \textsuperscript{(19)}.

%%K t03-c1-08-1 p
8. --- Accanto a lui c'è la \VUgras{lattaia} \textsuperscript{(24)} con il suo
\VUgras{asino} \textsuperscript{(23)}. Torna dalla città, dove ha portato il suo buon
latte per i bambini.

%%K t03-tit-5 sub
\VUsaut{4.0mm}
\VUpk{10.2pt}{40}{Il frutteto.}
\VUsaut{3.4mm}

%%K t03-c1-09-1 p
9. --- Messer Grigio (l'asino) è molto contento di questa corsa mattutina. Respira con
piacere l'aria profumata dai \VUgras{fiori} \textsuperscript{(20)} che coprono gli
\VUgras{alberi da frutto} \textsuperscript{(18)} del \VUgras{frutteto} vicino. Quei
\VUgras{peschi} \textsuperscript{(18)} e quegli \VUgras{albicocchi}
\textsuperscript{(19)} sono tutti rosa e bianchi, e promettono un buon raccolto.

%%K t03-c1-10-1 p
10. --- Intanto le piccole \VUgras{api} \textsuperscript{(22)} dell'\VUgras{arnia}
\textsuperscript{(21)} ci vanno a raccogliere il loro dolce \VUgras{miele},
ronzando mentre lavorano.

%%K t03-c1-11-1 p
11. --- Ma che cosa succede? Ecco i bambini del paese che arrivano di corsa e fanno
scappare le \VUgras{oche} \textsuperscript{(25)} spaventate. È il corteo che esce dalla
chiesa dopo il \VUgras{battesimo} del figlio del notaio.

%%K t03-c1-12-1 p
12. --- Il piccolo Giacomo, nella sua \VUgras{culla} \textsuperscript{(26)}, si sveglia
per lo scampanio allegro, e sua sorella Maddalena, che vuole vedere anche lei il
corteo, chiede alla mamma di venirla a pettinare e a vestire sulla soglia di casa.

%%K t03-c1-13-1 p
13. --- La mamma è d'accordo. Si mette il \VUgras{fazzoletto} \textsuperscript{(28)},
il suo \VUgras{corpetto} \textsuperscript{(29)} e il suo \VUgras{grembiule}
\textsuperscript{(30)}, e viene a sedersi su una seggiolina davanti alla porta.
Maddalena ha addosso solo la \VUgras{sottana} \textsuperscript{(31)} e il suo
\VUgras{corpetto} \textsuperscript{(32)}.

%%K t03-c1-14-1 p
14. --- Com'è contenta! Si dimentica dei suoi fiori preferiti, i suoi
\VUgras{garofani} \textsuperscript{(34)}, su cui si posano le \VUgras{farfalle}
\textsuperscript{(35)}, i suoi \VUgras{tulipani} \textsuperscript{(36)}, i suoi
\VUgras{mughetti} \textsuperscript{(37)} nei \VUgras{vasi} \textsuperscript{(38)} sul
\VUgras{muretto} \textsuperscript{(33)}, e le sue \VUgras{rose}
\textsuperscript{(39)} che fanno una siepe così bella. Guarda solo i bei vestiti delle
signore del corteo.

%%K t03-c1-15-1 p
15. --- I ragazzi del paese badano poco ai vestiti. Si buttano avanti e si spingono
per prendere i \VUgras{confetti} \textsuperscript{(55)} o le \VUgras{monete}
\textsuperscript{(52)}. Uno di loro piange \textsuperscript{(40)}.

%%K t03-c1-16-1 p
16. --- Un altro ha pestato la zampa del \VUgras{cane} \textsuperscript{(42)}, che
scappa guaendo e fa scappare la \VUgras{gallina} \textsuperscript{(43)} con i suoi
\VUgras{pulcini} \textsuperscript{(44)}.

%%K t03-tit-6 sub
\VUsaut{5.0mm}
\VUpk{10.2pt}{40}{Il corteo del battesimo.}
\VUsaut{4.0mm}

%%K t03-c1-17-1 p
17. --- In testa al corteo cammina la \VUgras{balia} \textsuperscript{(45)}. Ha una
bella \VUgras{cuffia} \textsuperscript{(46)} con lunghi nastri. Porta il
\VUgras{neonato} \textsuperscript{(47)}, tutto vestito di \VUgras{pizzi}
\textsuperscript{(48)}.

%%K t03-c1-18-1 p
18. --- Dietro la balia vengono il \VUgras{padrino} \textsuperscript{(49)} e la
\VUgras{madrina} \textsuperscript{(53)}. Distribuiscono i confetti e le monete. Il
padrino ha i \VUgras{guanti} \textsuperscript{(51)} e un \VUgras{cilindro}
\textsuperscript{(50)}. La madrina tiene in mano un bel \VUgras{mazzo di fiori}
\textsuperscript{(54)}.

%%K t03-c1-19-1 p
19. --- Dietro di loro vediamo lo \VUgras{zio} \textsuperscript{(56)} e la
\VUgras{zia} \textsuperscript{(58)} del bambino. La zia ha un elegante
\VUgras{cappello} \textsuperscript{(59)}, lo zio una \VUgras{redingote}
\textsuperscript{(57)} nera e un panciotto bianco. Poi vengono il \VUgras{cugino}
\textsuperscript{(60)} e sua \VUgras{cugina} \textsuperscript{(61)}. Lei tiene per mano
la \VUgras{sorella} \textsuperscript{(62)} del neonato, la piccola Berta.

%%K t03-c1-20-1 p
20. --- L'altro \VUgras{fratello} \textsuperscript{(63)} di Berta cammina dietro. Dà la
mano a un \VUgras{parente} \textsuperscript{(64)}. Gli \VUgras{amici}
\textsuperscript{(65)} della famiglia vengono dopo di loro.

%%K t03-tit-7 sub
\VUsaut{4.6mm}
\VUpk{10.2pt}{40}{Gli spettatori.}
\VUsaut{3.4mm}

%%K t03-c1-21-1 p
21. --- Vicino al corteo, un \VUgras{mendicante} \textsuperscript{(66)} è appoggiato
alla sua \VUgras{stampella} \textsuperscript{(67)}. Chiede l'elemosina.

%%K t03-c1-22-1 p
22. --- Dall'altra parte c'è un bambino molto \VUgras{educato}
\textsuperscript{(68)}. Si toglie il berretto e dà il \VUgras{«buongiorno»} alle
persone che conosce.

%%K t03-c1-23-1 p
23. --- Gli abitanti del paese sono usciti di casa per vedere passare il corteo del
battesimo. Vediamo un \VUgras{contadino} \textsuperscript{(69)} con il suo
\VUgras{berretto di cotone} e accanto a lui una \VUgras{contadina}
\textsuperscript{(70)}. Poi una \VUgras{vecchia} \textsuperscript{(71)} appoggiata al
suo \VUgras{bastone} \textsuperscript{(72)}. Infine il \VUgras{mugnaio}
\textsuperscript{(73)} con il suo largo \VUgras{cappello di feltro}
\textsuperscript{(74)}, e una giovane contadina con gli \VUgras{zoccoli}
\textsuperscript{(75)} che insegna al suo bambino a camminare.

%%K t03-c1-24-1 p
24. --- È una scena molto vivace.
\VUsaut{4.0mm}

%%K t03-tit-8 sub
\VUcentre{12.0pt}{0}{\textit{Seconda scena.}}
\VUsaut{3.0mm}

%%K t03-tit-9 sub
\VUpk{10.2pt}{40}{Un giorno di festa.}
\VUsaut{2.4mm}

%%K t03-tit-10 sub
\VUpk{10.2pt}{40}{Sport acquatici e regata.}
\VUsaut{3.0mm}

%%K t03-c2-01-1 p
1. --- L'\VUgras{estate} è arrivata. Il sole ardente del mese di \VUgras{giugno}
(luglio o \VUgras{agosto}) invita i giovani a fare il bagno e ad andare in barca.

%%K t03-c2-02-1 p
2. --- Sulla riva destra del fiume i rematori si spogliano per prendere parte agli
sport acquatici e alla regata.

%%K t03-c2-03-1 p
3. --- Uno di loro si toglie le \VUgras{scarpe} \textsuperscript{(1)}; le slaccia (le
sbottona). I suoi due amici sono già pronti. Ora hanno addosso solo la
\VUgras{maglia} \textsuperscript{(2)} e i \VUgras{calzoncini da bagno}
\textsuperscript{(3)}. Chiacchierano aspettando gli amici. Parlano della fiera e della
festa che si tiene sull'altra riva.

%%K t03-c2-04-1 p
4. --- Per terra, vicino a loro, ci sono i \VUgras{vestiti} dei loro amici: un
\VUgras{paio} di \VUgras{stivali} \textsuperscript{(4)}, \VUgras{scarpe}
\textsuperscript{(5)}, \VUgras{calze} \textsuperscript{(6)}, cappelli di
\VUgras{feltro} \textsuperscript{(7)} e di \VUgras{paglia} \textsuperscript{(8)} e una
\VUgras{cravatta} \textsuperscript{(9)}.

%%K t03-c2-05-1 p
5. --- Uno dei giovani ha piegato con cura la sua \VUgras{giacca}
\textsuperscript{(10)}. La posa su un asciugamano, nel quale il suo vicino avvolgerà
fra poco tutti e due gli \VUgras{abiti}.

%%K t03-c2-06-1 p
6. --- Un sesto ragazzo è in ritardo. Ha ancora il \VUgras{berretto}
\textsuperscript{(11)} in testa. Non si è tolto né la \VUgras{camicia}
\textsuperscript{(12)}, né il \VUgras{colletto} \textsuperscript{(13)}, né i
\VUgras{polsini} \textsuperscript{(14)}. Tiene in mano il \VUgras{panciotto}
\textsuperscript{(17)} e ha ancora i \VUgras{pantaloni} \textsuperscript{(16)} e le
\VUgras{bretelle} \textsuperscript{(15)}. Per la prossima gara non sarà certo pronto.

%%K t03-c2-07-1 p
7. --- Vicino ai giovani, un sentiero fiancheggiato di \VUgras{cardi}
\textsuperscript{(18)} scende verso il fiume, dove il vecchio Giacomo prepara, in
mezzo alle \VUgras{canne} \textsuperscript{(19)}, i \VUgras{fuochi d'artificio}
\textsuperscript{(20)} che si accenderanno la sera.

%%K t03-tit-11 sub
\VUsaut{4.4mm}
\VUpk{10.2pt}{40}{La gara.}
\VUsaut{3.4mm}

%%K t03-c2-08-1 p
8. --- Sul fiume, alcune signore al riparo dei loro parasoli sono in una \VUgras{barca
a remi} \textsuperscript{(21)}. Seguono la gara, che è molto emozionante. In ogni
\VUgras{yole} \textsuperscript{(22)} quattro \VUgras{rematori} \textsuperscript{(24)}
tirano con tutte le loro forze, mentre il \VUgras{timoniere} \textsuperscript{(26)}
tiene il \VUgras{timone} \textsuperscript{(25)} e guida la barca leggera. I
\VUgras{remi} \textsuperscript{(23)} battono l'acqua a tempo, e le barche leggere
volano a una velocità da capogiro.

%%K t03-c2-09-1 p
9. --- Alcuni giovani gareggiano, vicino al pilone del ponte, per il premio
dell'\VUgras{albero della cuccagna} \textsuperscript{(28)}. Uno di loro avanza con
prudenza sul palo elastico e scivoloso per arrivare alla \VUgras{bandierina}
\textsuperscript{(29)} fissata in cima. Il suo \VUgras{rivale} \textsuperscript{(30)},
meno fortunato, è caduto in acqua. Torna a riva a nuoto.

%%K t03-c2-10-1 p
10. --- I ragazzi più grandi fanno la \VUgras{giostra sull'acqua}
\textsuperscript{(32)} vicino alla \VUgras{banchina} \textsuperscript{(33)}. Una
grande \VUgras{folla} \textsuperscript{(31)} guarda tutti questi giochi, appoggiata al
\VUgras{parapetto} del \VUgras{ponte} \textsuperscript{(27)} e stipata lungo la
\VUgras{riva}. Qualche spettatore però si volta a seguire il gioco della
\VUgras{pertica} \textsuperscript{(45)}, o ad ammirare il \VUgras{corteo del sindaco}
che va alla \VUgras{premiazione}.

%%K t03-c2-11-1 p
11. --- Per primi ecco alcuni \VUgras{monelli} \textsuperscript{(35)} che corrono
davanti alla \VUgras{banda} \textsuperscript{(36)}; poi vengono il \VUgras{sindaco}
\textsuperscript{(37)}, i suoi assessori e il \VUgras{consiglio comunale} della
cittadina. Per ultimi sfilano i \VUgras{pompieri} \textsuperscript{(38)}.

%%K t03-tit-12 sub
\VUsaut{5.0mm}
\VUpk{10.2pt}{40}{La fiera.}
\VUsaut{3.6mm}

%%K t03-c2-12-1 p
12. --- Sulla \VUgras{piazza della fiera} \textsuperscript{(39)} c'è una gran folla. La
gente viene a vedere i vari \VUgras{baracconi}: la \VUgras{giostra}
\textsuperscript{(40)}, il \VUgras{circo} \textsuperscript{(41)}, dove lo spettacolo è
già cominciato; il \VUgras{ballo pubblico} \textsuperscript{(42)}, dove i giovani della
città si godono il piacere di \VUgras{ballare}.

%%K t03-c2-13-1 p
13. --- In fondo vediamo l'\VUgras{arena} \textsuperscript{(44)}, dove il coraggioso
\VUgras{matador} spagnolo combatte il \VUgras{toro} \textsuperscript{(45)}
infuriato; poi le \VUgras{montagne russe} \textsuperscript{(49)} e il \VUgras{tiro a
segno} \textsuperscript{(46)}, dove ci si esercita a maneggiare le \VUgras{armi da
fuoco} e a sparare al \VUgras{bersaglio}.

%%K t03-c2-14-1 p
14. --- Lì accanto c'è il \VUgras{teatro di fiera} \textsuperscript{(47)}, dove si
recitano la \VUgras{commedia} e il \VUgras{dramma}. Vicinissimo c'è il baraccone del
\VUgras{gigante} \textsuperscript{(48)} e del \VUgras{nano}. La \VUgras{statura} del
primo è di due metri e quindici; il secondo è alto appena quanto un bambino.

%%K t03-c2-15-1 p
15. --- Per ultimo ecco il grande \VUgras{serraglio} \textsuperscript{(50)} con i suoi
\VUgras{animali feroci}, la sua \VUgras{tigre} \textsuperscript{(51)} dagli
\VUgras{artigli} potenti, il suo \VUgras{leone} \textsuperscript{(52)} dalla folta
\VUgras{criniera}, le sue due \VUgras{giraffe} \textsuperscript{(53)} e il suo
\VUgras{elefante} \textsuperscript{(54)}, che usa la sua \VUgras{proboscide} con tanta
abilità.

%%K t03-c2-16-1 p
16. --- Il \VUgras{domatore} \textsuperscript{(55)} che addestra questi animali è in
piedi alla porta del baraccone.
\VUsaut{6.0mm}
\VUfilet{22mm}
